\section{情報を扱う責任：公正な条件と説明の相手}
\label{theory:information:fairness}
情報通信技術は、場所や時間を越えて情報へ接近する能力を増す。試験中の検索や外部への質問が可能でも、その使用が試験制度の目的に適うとは限らない。試験は一定の条件の下で能力を比較するための測定である。ある受験者だけが禁止された助力を受ければ、点数が示す意味が変わり、他の受験者との比較を損なう。

発覚による自分の不利益を引き受けるという判断だけでは、他人の選考機会への影響を取り除けない。技術倫理では、行為者に生じる便益と危険だけでなく、情報にアクセスできない人、評価される人、結果に基づいて判断する組織へ及ぶ影響を考える。刑法の保護する利益も身体や物だけではない。平成23年度問3の注記は当時の刑法233条を挙げるが、逮捕という記載から個別事件の有罪を推論しない。

\section{暗黙知と形式知：技能の共有と知識の再利用}
\label{theory:information:knowledge}
装置の異音を聞いて故障の前兆へ気付く技能は、手順を読んだだけではすぐに身に付かない。一方、測定の設定・判定の条件を文章や図で残せば、他の担当者が比較や改善を行いやすい。前者の経験・技能として保持される知を暗黙知、後者の明示的な表現を持つ知を形式知という。二つは相互に補完する\parencite{jaist_knowledge}。

熟練者と共に作業して勘所を学ぶこと、経験を条件付きの手順へ表すこと、複数の手順や記録を組み合わせること、手順を実際に試して技能へ取り込むことをそれぞれ共同化・表出化・連結化・内面化と整理できる。ただし技能を完全に文章化できると仮定するモデルではない。手順書に温度や振動の条件が書かれていても、異常時の判断には実地の訓練を要することがある。共有可能性と、知識が成立した文脈の保持を両立させる。

\section{機密性・完全性・可用性：何を失う危険か}
\label{theory:information:cia}
情報セキュリティは秘密を守ることだけではない。患者の記録が漏れなくても数値が改ざんされれば診療判断を誤る。正しい記録が保管されていても必要時に閲覧できなければ業務が止まる。機密性は権限のない開示を防ぐこと、完全性は情報や処理の正確さと完全な状態を保つこと、可用性は認可された人が必要時に利用できることを扱う。

\begin{center}
\begin{tabular}{@{}p{0.18\linewidth}p{0.35\linewidth}p{0.40\linewidth}@{}}
\toprule
要素 & 失敗の例 & 対策の目的\\
\midrule
機密性 & 顧客へ他の顧客の連絡先を送る & 役割に応じた権限と送信先確認\\
完全性 & 測定値や履歴が変更される & 正確な更新、変更記録、整合検査\\
可用性 & 必要時に記録へ接近できない & 運用継続、復旧、適切な冗長化\\
\bottomrule
\end{tabular}
\end{center}
即時性や保存性という日常語を、この三要素の標準的名称と取り違えない。三要素は同じ対策で同時に改善する場合も、緊急時のアクセス制限と利用可能性のように調整を要する場合もある。資産と業務の条件に応じて要求水準を定める。

\section{方針と管理体系：文書から継続する運用へ}
\label{theory:information:management}
情報セキュリティポリシーは組織の方針・行動指針を示す。情報セキュリティマネジメントシステム（ISMS）は、対象資産とリスクを把握し、方針を計画と実施へ結び付け、結果を評価して改善する体系である。方針を作っただけで、その実施が保証されるわけではない。幹部を含む教育、担当と権限の設定、監視、事故時の初動と復旧、原因分析を継続して行う。

平成23年度問10の原典は2005年のISO27000と記すが、公開書誌のISO/IEC27000初版は2009年である\parencite{iso27000_2009}。原典の名称を黙って別規格へ書き換えず、空欄のCIA定義の判断と書誌上の疑義を分ける。規格全条項はこの初稿で照合しておらず、一般的な改善策を各版の逐語的要求事項とは断定しない。

\section{個人情報の範囲：公開・評価・映像でも識別性を見る}
\label{theory:privacy:scope}
個人情報保護は、情報の利用による便益と、本人の権利利益への影響を調整する。氏名等の事実だけでなく、個人についての評価や、識別可能な音声・映像も対象となり得る。死者だけに関する情報は生存する個人の定義から外れるが、それが生存する家族の情報でもある場合は、その家族についての識別性を調べる。新聞掲載や公開サイトに載ったというだけでも個人情報性は消えない\parencite{privacy_general}。

組織内の従業員番号等は、番号から本人へ容易に照合できるかを確認する。意味のない乱数でも照合表を管理していれば、単なる数字だから除外するという判断はできない。一方、全ての識別不能な番号へ同じ結論を無条件に当てはめない。個人情報、検索可能な集合としてのデータベース、その構成要素である個人データという段階も区別する。

\subsection{旧小規模事業者の基準と改正}
\label{theory:privacy:history}
平成24年度問4の5000人基準は当時の事業者適用範囲についての問題である。個々のデータベースを5000人以下へ分割すれば事業者の規律を免れる、という判断は成立しない。2017年5月30日施行の改正で5000人以下を適用除外する規定は廃止された\parencite{privacy_small2017}。現在の制度の説明へ旧閾値を残さない。旧政令の計数の全条件と当時版の直接照合は査読へ残す。

\section{本人同意：利用目的と第三者提供を分ける}
\label{theory:privacy:consent}
情報を取得する際の目的、内部で使う目的、他の主体へ提供することは、同じ行為ではない。ポイント管理のために得た顧客電話番号を求人へ転用する場合には、取得目的に必要な範囲かを検討する。名簿を内部で作成できることから、その名簿を他の家庭へ配布できるとは推論しない。

本人同意は、示された取扱方法について承諾する意思表示である。口頭や書面等の明示だけに限られず、黙示の同意が成立するかは個別具体的な事実で判断する\parencite{privacy_consent2017,privacy_consent_current}。したがって、問い合わせを受けたこと、説明の便宜があること、本人に役立つと事業者が考えることだけで、第三者へ連絡先を渡す承諾を補わない。

令和3年度問14では、カード請求の調査依頼から加盟店への連絡先提供まで黙示同意があると初回に解したが、公式正答は全記述不適切である。提供方法への承諾が確認できる契約・発言等が設問にない点を記録し、初回解答を保持して査読対象とする。黙示同意が法的に一切存在しない、と単純化して公式へ合わせない。CC一斉送信も、宛先同士へアドレスが見えるという実際の情報流れから判断する。

\subsection{配送：宛名情報と封入物の中身}
\label{theory:privacy:delivery}
発送のために氏名・住所を業者へ伝える場合、その宛名情報の取扱いは委託となり得る。他方、封入された媒体の中身を通常の配送業者は取り扱わない。中身の取扱合意等がなければ、その中身の個人データの取扱いを委託したとは一般に考えない\parencite{privacy_delivery}。この区別は、発送者の安全管理義務や配送契約上の責任を全て消すものではない。安全な配送方法、適切な業者、暗号化等を情報の条件に応じて検討する。

\subsection{災害時：必要性と提供範囲を限定する}
\label{theory:privacy:evacuation}
2013年の災害対策基本法改正は避難行動要支援者名簿の作成を市町村へ義務付けた。平時の支援者への名簿提供と、発災又は発生のおそれのある場合の提供は条件が違う\parencite{privacy_evacuation2013,privacy_evacuation_scope}。災害時には生命身体の保護へ特に必要な場合、必要な限度で同意なし提供が認められる。個人情報だから一切提供できない、又は災害対策という目的なら全情報を自由に共有できる、のどちらにも一般化しない。

\section{個人識別符号：生体特徴の列挙と変換条件}
\label{theory:privacy:biometrics}
政令はDNA、顔、虹彩、声、歩行、静脈、指紋・掌紋の七つの生体特徴を列挙する。電子計算機向けに変換され、特定個人の識別に足りる基準に合う符号という条件がある\parencite{privacy_general}。生体の特徴全てがそのまま個人識別符号になるわけではない。令和元年度問4は七分類が全て含まれること、令和6年度問9は骨振動の特性という非列挙項目を見分けることを問う。認証技術として実現可能かという一般論を、法令が何を列挙しているかという設問へ置換しない。

\section{説明責任と守秘：何を誰へ伝えるか}
\label{theory:information:accountability}
説明責任では、判断に影響した条件、根拠、限界、利益相反を示す。研究の実施場所や領域が違っても、その価値は共有される\parencite{singapore2010}。相手が質問し、理由を吟味する双方向の対話を支えることと、相手へ無条件の納得を要求することは違う。守秘を要する場合には、情報を抽象化する、説明範囲を確認する、権限のある監督者へ相談するなどにより、説明と秘密保持の双方を満たす方法を考える。
